This thesis set out to provide a reproducible, release-wise account of how the persistent packages of Debian's minbase installation have developed across the distribution's full stable history, a longitudinal characterisation that the existing literature had not assembled. To that end it built a containerised analysis pipeline and applied it to the 25 source packages of the persistent core, those present in the minbase installation of at least 11 of the 13 stable releases from Potato (2000) to Trixie (2025), extracting five metrics per release: installed size, dependency structure, lines of code, per-function cyclomatic complexity, and distribution-specific patch volume.

The resulting artefacts describe a base system that grew in size while remaining stable in the complexity of its parts. The persistent core grew roughly 5.3-fold in lines of code and roughly fourfold in installed footprint over the 25-year period, yet the median per-function cyclomatic complexity held at 3 in every release and the median per-package dependency count at 1 in every release after Potato. The core therefore expanded through the addition of many small, low-complexity, loosely coupled units rather than through an increase in the complexity of the typical unit. The composition of the minbase changed non-monotonically across releases, and its largest movements, the \texttt{systemd} init-system transition at Jessie and the recent decline in package count, are attributable to packaging structure rather than to changes in capability. The volume of distribution-specific patches, comparable only from the adoption of the \texttt{3.0 (quilt)} source format at Jessie onward, declined across the most recent releases and was dominated in every release by \texttt{glibc}; the file-level structure of those patches changed in parallel, from a single bundled \texttt{.diff.gz} per package to many named, purpose-specific patch files.

Alongside these descriptive results, the thesis delivers the means to reproduce and extend them. The pipeline, built from established open-source tools (\texttt{debootstrap}, \texttt{tokei}, \texttt{pmccabe}, and \texttt{gnuplot}) and orchestrated with Docker and a Makefile, regenerates the full set of per-release CSV files and figures from the Debian archives. An interactive companion application, the Debian Bloat Explorer, reads the same CSV files and supports drill-down into individual packages and releases; it was the instrument through which several measurement artefacts, and the patch-metric inconsistency at the source-format transition, were diagnosed.

The contribution of this work is a descriptive baseline rather than an explanatory account. The artefacts characterise what has happened to Debian's persistent core; they do not establish why, and the design supports no causal or correlational claim relating one dimension to another. Whether the divergence between rising size and stationary per-function complexity reflects a general property of long-lived distributions, and what mechanisms produce it, are questions that the reproducible pipeline and per-release data presented here provide a basis to investigate in future work.
